\chapter{Solutions to the Project Derivation Exercises}
\label{app:solutions}

\begin{pitfall}[read this first]
These are the Jacobians that \texttt{AGENTS.md} (Guardrail 4) asks you to derive yourself. Attempt each exercise first, check your
result numerically, and only then compare. Every result below was verified against central finite differences built on the project's
$\bplus/\bminus$ operators.
\end{pitfall}

% ------------------------------------------------------------------------------
\section{Error-state dynamics \texorpdfstring{$F_c$}{Fc} and \texorpdfstring{$G_c$}{Gc}}
\label{sol:Fc}

\paragraph{Attitude.} Differentiate $R = \hat R\Exp(\delta\theta)$ with \cref{lem:dexp}, to first order:
\begin{equation*}
  \dot R = \dot{\hat R}\Exp(\delta\theta) + \hat R\Exp(\delta\theta)\sx{\dot{\delta\theta}}
         = \hat R\sx{\hat\omega}\Exp(\delta\theta) + \hat R\Exp(\delta\theta)\sx{\dot{\delta\theta}} .
\end{equation*}
The true kinematics give $\dot R = R\sx{\omega} = \hat R\Exp(\delta\theta)\sx{\omega}$. Equate and multiply on the left by $(\hat R\Exp(\delta\theta))\T$:
\begin{equation*}
  \sx{\omega} = \Exp(-\delta\theta)\sx{\hat\omega}\Exp(\delta\theta) + \sx{\dot{\delta\theta}} = \sx{\Exp(-\delta\theta)\hat\omega} + \sx{\dot{\delta\theta}} .
\end{equation*}
With $\Exp(-\delta\theta)\hat\omega \approx \hat\omega - \delta\theta\times\hat\omega = \hat\omega + \hat\omega\times\delta\theta$ and $\omega - \hat\omega = -\delta b_g - n_g$:
\begin{equation}
  \dot{\delta\theta} = -\sx{\hat\omega}\,\delta\theta - \delta b_g - n_g .
\end{equation}
The term $-\sx{\hat\omega}\delta\theta$, absent in the global version (\cref{ex:global-attitude}), appears because the error is expressed in body axes
that rotate with angular velocity $\hat\omega$. An error fixed in the world appears to rotate backwards in the body frame.

\paragraph{Velocity.} With $R \approx \hat R(I + \sx{\delta\theta})$ and $a_m - b_a - n_a = \hat a - \delta b_a - n_a$:
\begin{equation*}
  \dot{\hat v} + \dot{\delta v} \approx \hat R\hat a + \hat R\sx{\delta\theta}\hat a - \hat R\delta b_a - \hat Rn_a + g_w,
\end{equation*}
dropping the second-order products $\sx{\delta\theta}\delta b_a$ and $\sx{\delta\theta}n_a$. Subtracting $\dot{\hat v} = \hat R\hat a + g_w$ and using
$\sx{\delta\theta}\hat a = -\sx{\hat a}\delta\theta$:
\begin{equation}
  \dot{\delta v} = -\hat R\sx{\hat a}\,\delta\theta - \hat R\,\delta b_a - \hat R\,n_a .
\end{equation}

\paragraph{Assembly.} With the frozen orders $\delta x = [\delta p, \delta v, \delta\theta, \delta b_a, \delta b_g]$ and $w = [n_a, n_g, n_{b_a}, n_{b_g}]$:
\begin{equation}
  F_c = \begin{bmatrix}
    0 & I & 0 & 0 & 0 \\
    0 & 0 & -\hat R\sx{\hat a} & -\hat R & 0 \\
    0 & 0 & -\sx{\hat\omega} & 0 & -I \\
    0 & 0 & 0 & 0 & 0 \\
    0 & 0 & 0 & 0 & 0
  \end{bmatrix}, \qquad
  G_c = \begin{bmatrix}
    0 & 0 & 0 & 0 \\
    -\hat R & 0 & 0 & 0 \\
    0 & -I & 0 & 0 \\
    0 & 0 & I & 0 \\
    0 & 0 & 0 & I
  \end{bmatrix},
\end{equation}
all blocks $3\times3$, with $\hat R = \hat R_{wb}$, $\hat a = a_m - \hat b_a$, $\hat\omega = \omega_m - \hat b_g$. The structural zeros follow from the cascade: position
depends only on velocity; attitude does not depend on position, velocity or the accelerometer bias; the biases are driven only by
their own noise. Since $Q_c$ is block-diagonal and $\hat R\hat R\T = I$,
$Q_d \approx G_cQ_cG_c\T\Dt = \diag(0,\ \sigma_a^2I,\ \sigma_g^2I,\ \sigma_{b_a}^2I,\ \sigma_{b_g}^2I)\,\Dt$.

\paragraph{Checking the signs.} Rotating a positive error $\delta\theta$ about $\hat\omega$ itself gives $-\sx{\hat\omega}\hat\omega = 0$: an error about the rotation axis
does not move, as it should. A positive gyro bias error makes the true rate smaller than the estimated one, so the true attitude
lags and $\dot{\delta\theta} = -\delta b_g$. For a level, static IMU, $\hat a = [0, 0, g]$, and $-\sx{\hat a}\delta\theta = g\,(\delta\theta_y, -\delta\theta_x, 0)$: a tilt about $y$
produces a velocity error along $x$, the gravity leak.

% ------------------------------------------------------------------------------
\section{Closed-form blocks of \texorpdfstring{$\Phi$}{Phi}}
\label{sol:Phi}

The project's default is $\Phi \approx I + F_c\Dt$. If a higher-order transition is adopted (and logged), the exact solution of the error
dynamics over a step with constant $\hat\omega$, $\hat a$, and with $\hat R(s) = \hat R_k\Exp(\hat\omega s)$ evolving inside the step, has these blocks,
with $\phi = \hat\omega\Dt$ and $N(\phi)$ from \cref{prop:exact-prop}:
\begin{align*}
  \Phi_{\theta\theta} &= \Exp(-\phi) = \Exp(\phi)\T, & \Phi_{\theta b_g} &= -\Dt\,\Jr(\phi), \\
  \Phi_{v\theta} &= -\hat R_k\sx{\Dt\,\Jl(\phi)\hat a}, & \Phi_{vb_a} &= -\hat R_k\,\Dt\,\Jl(\phi), \\
  \Phi_{p\theta} &= -\hat R_k\sx{\Dt^2N(\phi)\hat a}, & \Phi_{pb_a} &= -\hat R_k\,\Dt^2N(\phi), \qquad \Phi_{pv} = \Dt\,I .
\end{align*}
The remaining coupling blocks ($\Phi_{vb_g}$, $\Phi_{pb_g}$) have no compact closed form; compute them by series or numerically, and use
Van Loan's method (\cref{prop:vanloan}) or a fine numerical integration of $\dot\Phi = F_c(s)\Phi$ as the test oracle. To first order in
$\Dt$, every block reduces to the corresponding block of $I + F_c\Dt$.

% ------------------------------------------------------------------------------
\section{Pose measurement Jacobian}
\label{sol:Hpose}

With $h(x) = (p_{wb}, q_{wb})$ and $r(\delta x) = h(\hat x\bplus\delta x)\bminus h(\hat x)$:
\begin{equation*}
  r_p = (\hat p + \delta p) - \hat p = \delta p, \qquad
  r_\theta = \Log_q\big(\hat q\inv\qprod\hat q\qprod\Exp_q(\delta\theta)\big) = \delta\theta .
\end{equation*}
Both are exact, so
\begin{equation}
  H = \begin{bmatrix} I_3 & 0 & 0 & 0 & 0 \\ 0 & 0 & I_3 & 0 & 0 \end{bmatrix} \in \R^{6\times15}.
\end{equation}
Its nullspace is spanned by $\delta v$, $\delta b_a$ and $\delta b_g$: these are not measured directly and are corrected through their
cross-covariances with $\delta p$ and $\delta\theta$. The simplicity of the attitude block is a consequence of the noise model
$q_{meas} = q\qprod\Exp_q(n_\theta)$ matching the local perturbation, which makes $r_\theta = \Log(\Exp(\delta\theta)\Exp(n_\theta)) \approx \delta\theta + n_\theta$.
If the residual is instead linearized around a non-zero value $r_0$ (an iterated update, or the optimization view of
\cref{sec:manifold-update}), the derivative of $\Log(\Exp(-\delta\theta)\Exp(r_0))$ with respect to $\delta\theta$ is $-\Jlinv(r_0)$, which reduces to $-I$ for small
$r_0$, the sign being that of the optimization convention.

% ------------------------------------------------------------------------------
\section{Cloning Jacobian}
\label{sol:clone}

At cloning time the true clone equals the true IMU pose and the clone's nominal equals the IMU nominal, so
\begin{equation*}
  \delta p_c = p - \hat p = \delta p, \qquad \Exp(\delta\theta_c) = \hat q\inv\qprod q = \Exp(\delta\theta)\ \Rightarrow\ \delta\theta_c = \delta\theta,
\end{equation*}
exactly, without linearization. With the new clone appended at the end of an $n = 15 + 6N$ dimensional state,
\begin{equation}
  J_g = \begin{bmatrix} I_3 & 0 & 0 & 0 & 0 & 0_{3\times6N} \\ 0 & 0 & I_3 & 0 & 0 & 0_{3\times6N} \end{bmatrix}\in\R^{6\times n},
  \qquad J = \begin{bmatrix} I_n \\ J_g \end{bmatrix}.
\end{equation}
Consequently $JPJ\T$ copies blocks: the new clone's covariance is
$\begin{bmatrix} P_{pp} & P_{p\theta} \\ P_{\theta p} & P_{\theta\theta}\end{bmatrix}$, and its cross-covariance with every other state is the stacked $p$ and $\theta$ rows of $P$.
No matrix product is needed in the implementation.

If the clone stored the camera pose, $R_{wc} = R_{wb}R_{bc}$ and $p_{wc} = p_{wb} + R_{wb}p_{bc}$ would give
$\delta\theta_c = R_{bc}\T\delta\theta$ and $\delta p_c = \delta p - \hat R\sx{p_{bc}}\delta\theta$, by the composition and rotation-action rows of \cref{tab:jacobians}. Storing the
IMU pose keeps $J$ trivial and moves the extrinsics into the measurement model.

% ------------------------------------------------------------------------------
\section{Visual measurement Jacobians}
\label{sol:Hcam}

Write $\hat d_i = \hat R_{wb_i}\T(\hat p_f - \hat p_{wb_i})$, the feature in the clone's body frame. Perturbing the clone and the feature,
\begin{equation*}
  R_{wb_i}\T(p_f - p_{wb_i}) \approx (I - \sx{\delta\theta_i})\hat R_{wb_i}\T(\hat p_f + \delta p_f - \hat p_{wb_i} - \delta p_i)
  \approx \hat d_i + \sx{\hat d_i}\delta\theta_i - \hat R_{wb_i}\T\delta p_i + \hat R_{wb_i}\T\delta p_f ,
\end{equation*}
using the inverse-action row of \cref{tab:jacobians}. Applying $R_{bc}\T$ (the subtraction of $p_{bc}$ is constant):
\begin{equation}
  \frac{\partial p_{c_i}}{\partial\delta\theta_i} = R_{bc}\T\sx{\hat d_i}, \qquad
  \frac{\partial p_{c_i}}{\partial\delta p_i} = -R_{bc}\T\hat R_{wb_i}\T, \qquad
  \frac{\partial p_{c_i}}{\partial p_f} = R_{bc}\T\hat R_{wb_i}\T .
\end{equation}
Moving the clone by $\delta$ is equivalent to moving the feature by $-\delta$, hence the opposite signs. With
$J_\pi = \frac1{\hat Z}\begin{bmatrix} I_2 & -\pi(\hat p_{c_i})\end{bmatrix}$ from \cref{prop:proj-jacobian}:
\begin{equation}
  H_{x,ji} = \Big[\ 0_{2\times15}\ \ \cdots\ \ \underbrace{-J_\pi R_{bc}\T\hat R_{wb_i}\T}_{\delta p_i}\ \ \underbrace{J_\pi R_{bc}\T\sx{\hat d_i}}_{\delta\theta_i}\ \ \cdots\ \Big],
  \qquad
  H_{f,ji} = J_\pi R_{bc}\T\hat R_{wb_i}\T ,
\end{equation}
with the two non-zero blocks in clone $i$'s columns $15 + 6i \dots 20 + 6i$, and zeros in the IMU columns and in every other clone's columns.
The stacked $H_{x,j}$ therefore has one $2\times6$ non-zero block per observation (\cref{fig:hx-sparsity}).

\paragraph{Checking the signs.} Translate the clone along the camera's optical axis toward the feature: the feature gets closer, so its
image moves away from the principal point, and $-J_\pi R_{bc}\T\hat R_{wb_i}\T\delta p_i$ predicts exactly that. Rotate the clone about the camera's
$y$-axis by a small positive angle ($\delta\theta_i = \alpha R_{bc}e_y$): the image of a distant point on the optical axis moves in $-x$, which the
$\delta\theta_i$ block reproduces (for a near point, the lever arm $p_{bc}$ adds a small translation of the camera center).
